% !TeX program = pdflatex
% !TeX encoding = UTF-8
% Preamble.
\documentclass[12pt]{article}

% Packages.
\usepackage[utf8]{inputenc}     % Unicode support (Umlauts etc.)
\usepackage{amsmath, amssymb}   % Advanced math typesetting
\usepackage{stmaryrd, bbm}
\usepackage[french]{babel}      % Change hyphenation rules
\usepackage[T1]{fontenc}
\usepackage{geometry}
\usepackage{listings}           % Source code formatting and highlighting
\usepackage{csquotes}
\usepackage{interval}
\usepackage{pifont}

% Pages.
\sloppy 
\geometry{a4paper, total={17cm, 25cm}}
\renewcommand{\baselinestretch}{1.4}

% Macros.
%% Double interval.
\newcommand{\iinterval}[2]{
    \left\llbracket 
        #1; #2 
    \right\rrbracket
}
%% Finite set, with curly braces.
\newcommand{\setf}[1]{\left\{#1\right\}}
%% Parenthesis.
\newcommand{\of}[1]{\left(#1\right)}
%% Square brackets.
\newcommand{\brackets}[1]{\left[#1\right]}
%% Floor brackets.
\newcommand{\floor}[1]{\left\lfloor #1 \right\rfloor}
%% Ceiling brackets.
\newcommand{\ceil}[1]{\left\lceil #1 \right\rceil}
%% Absolute value.
\newcommand{\abs}[1]{\left\lvert #1 \right\rvert}
%% Pre-formatted "Exemple".
\newcommand{\Ex}{\textsc{Exemple}: }
%% Pre-formatted "Tel que".
\newcommand{\tq}{\text{ tel que }}
%% Items.
\newcommand{\In}{$\sim$}
\newcommand{\Il}{$\square$}
\newcommand{\Ia}{$\blacksquare$}
\newcommand{\Ii}{\ding{93}}

\AtBeginDocument{\def\labelitemi{\Ia}}

% Data.
\title{Terminale - Notations mathématiques.}
\author{Hector GARLATTI}
\date{\today{}}

\begin{document}
\maketitle

\section{Utilité.}
Ce document est une fiche, rédigée en \LaTeX, d'élève de fin de Terminale. 
Surtout destinée aux révisions pour les classes préparatoires, 
elle peut servir d'aperçu pour les élèves de Première et de lecture pour
le plaisir des mathématiques.

\subsection{Introduction.}
Le langage mathématique permet une expression objective, non-ambiguë
des théorèmes et des propriétés. Il peut est déjà utile en début
de Terminale de se familiariser avec ses notions car elles permettent:
\begin{itemize}
    \item la précision et la rigueur;
    \item de commencer à s'immerger dans l'univers
    des "mathématiques sérieuses";
    \item exprimer des concepts plus avancés.
\end{itemize}

Cependant, ces notations ne sont pas de excuses pour ne pas rédiger
et sont d'abord des outils pour enrichir la rédaction, non la simplifier.

\subsection{Raccourcis de cette fiche.}
Les notions et notations sont classées par leur icône:
\begin{itemize}
    \item[\In] Personnelle ou non conventionnelle.
    \item[\Il] Hors programmes ou moins critique.
    \item[\Ia] Classique, usuelle et utile.
    \item[\Ii] Clef, critique. 
\end{itemize}

\subsection{Règles pour la démonstration.}
En fin de Terminale, certains règles précises doivent être maîtrisées. 
En voici quelques unes:
\begin{enumerate}
    \item Un discours mathématiques n'est pas un suite de symboles.
    \item Tout objet est \textit{déclaré} avant d'être utilisé. \\
        \Ex \enquote{\underline{Soit} $x$ un nombre réel positif.}
    \item Une bonne démonstration est structurée par des locutions 
    qui guident le lecteur, que nous pouvons mettre en évidence
    (souligné, italique). \\
        \Ex \enquote{
            \underline{Montrons que} ... 
            \underline{Ainsi} ...
        }
\end{enumerate}


\section{Notations, quantificateurs.}
Ici ne seront que listé et brièvement mis en contexte les notation.
Pour des explications plus approfondies, se référer au leçons.

\subsection{Ensembles.}

\paragraph{Ensembles de nombres usuels.}
\begin{itemize}
    \item $\varnothing = \setf{}$: l'ensemble vide.
    \item[\Ii] $\mathbb{N}$: les entiers naturels, \\
        \Ex $\setf{0; 1; 2; \dots}$. 
    \item[\Ii] $\mathbb{Z}$: les entiers relatifs, \\ 
        \Ex $\setf{0; 1; -1; 2; -2; ...}$.
    \item $\mathbb{Q}$: les nombres rationnels, toutes les fractions: 
    $$ \setf{
        \frac{p}{q} \; | \;    
        \forall p; q \in \mathbb{Z}\times \mathbb{N}^{\star} 
    } $$
    \item[\Ii] $\mathbb{R}$: le continuum des réels, \\
        \Ex $\setf{\sqrt{2}; \pi; \sin(2); \dots}$.
    \item[\Ii] $\mathbb{C}$: les complexes, \\
        \Ex $\setf{1 + 2i; 2e^{\frac{1}{4}i\pi}; \dots}$.
    \item $\mathbb{R}^{n}$: les vecteurs (vecteurs, points, $\dots$) de $n$ dimensions. \\
        \Ex $(1; 2) \in \mathbb{R}^2$
    \item[\Il] $E[X]$: les polynômes à valeurs dans $E$. \\
        \Ex $(x^2 + 3) \in \mathbb{R}\left[X\right]$ 
    \item[\Il] $\mathcal{M}_{n, m}(E)$: les matrices de taille $(n, m)$ 
        à coefficients dans $E$ \\
        \Ex $\mathbbm{1}_3 \in \mathcal{M}_{3}(\mathbb{N})$
\end{itemize}

\paragraph{Segments et \enquote{regroupement} de nombres.} 
\begin{itemize}
    \item $\interval{a}{b}, \interval[open]{a}{b}$: 
        un intervalle ou \textit{segment} de $\mathbb{R}$.
    \item[\Il] $\iinterval{a}{b}$: un intervalle d'entier. \\
        \Ex $\iinterval{2}{5} = \setf{2; 3; 4; 5}$

    \item $T = \underbrace{\left(a; b; c; ...\right)}_{t}$ 
        est un $t$-uplet:
        \begin{itemize}
            \item son nombre d'éléments: $t = \text{Card}({E}) = |E|$, \\
                pour $t = 2$, $T$ est un \text{couple};
                pour $t = 3$, $T$ est un \text{triplet}, $\dots$;
            \item l'ordre compte; \\
            \Ex $\left(1; 2\right) \neq \left(2; 1\right)$
            \item les répétitions sont autorisées; \\
            \Ex $\left(1; 2; 2\right) \neq \left(1; 2\right)$
        \end{itemize} 

    \item $E = \setf{a; b; c; ...}$ est un ensemble:
        \begin{itemize}
            \item son nombre d'éléments: $\text{Card}({E}) = |E|$, \\
                pour $|E| = 1$, $E$ est un \textit{singleton};
            \item l'ordre ne compte pas; \\
            \Ex $\setf{1; 2} = \setf{2; 1}$
            \item les répétitions sont "effacées"; \\
            \Ex $\setf{1; 2; 2} = \setf{1; 2}$
        \end{itemize} 

    \item Création d'ensemble \enquote{par compréhension},
        de sous-ensemble de $E$ selon une structure: $$
        \setf{f(x) \;\vert\;  \forall x \in E} = f(E) 
        $$ 
        \Ex Tout les nombres impairs: $$
            \setf{2k + 1 \;\vert\; \forall k \in \mathbb{Z}} 
          = 2\mathbb{Z} + 1
        $$
\end{itemize}
\textsc{Attention}: Guillemets car les \textit{classes} et les 
    \textit{groupes} sont encore d'autres notions !


\paragraph{Appartenance.} 
Notations permettant de déclarer un object (\textit{cf.} \textbf{1.2.3})
de le décrire.

\begin{itemize}
    \item[\Ii] $x \in E$: un scalaire, un nombre, un vecteur fait parti d'un ensemble $E$. \\
        \Ex 
        \enquote{$x \in \mathbb{R}$} signifie 
        \enquote{\underline{Soit} $x$ un nombre réel.}
    \item $A \subset B$: un ensemble $A$ est compris dans, est un sous-ensemble de $B$. \\
        \Ex $A \subset B$ \\
        Ainsi, nous pouvons écrire pour nos ensembles usuels 
        (penser au diagramme des cercles): 
        \begin{equation}
            \varnothing \subset \mathbb{N} \subset \mathbb{Z} \subset 
            \mathbb{Q} \subset \mathbb{R} \subset \mathbb{C}
        \end{equation}
\end{itemize}

\paragraph{Altérations et operations.} 
Pour un ensemble $E$.
\begin{itemize}
    \item $E^{\star}$: excluant $0$. \\
        \Ex $\mathbb{N}^{\star}$, tout les entiers naturel 
        \textit{hormis} $0$.
    \item $E_{+}$: gardant uniquement les positifs. \\
        \Ex $\mathbb{R}_{+}$, tout les réels \textit{positifs}.
    \item $E_{-}$: gardant uniquement les négatifs. \\
        \Ex $\mathbb{R}^{\star}_{-}$, tout les réels négatifs sans $0$.
    
    \item[\Il] $A \setminus B$: difference d'ensembles. \\
        \Ex $\mathbb{R} \setminus \{1\}$, tout les réels excluant $1$.
    \item $A \cup B$: union 
        (ensembles, probabilités, \enquote{ou inclusif} logique, géométrie, $\dots$).\\
        \Ex $\{1; 2\} \cup \{2; 3\} = \{1; 2; 3\}$
    \item $A \cap B$: intersection 
        (ensembles, probabilités, \enquote{et} logique, géométrie, $\dots$).\\
        \Ex \begin{itemize}
            \item $\{1; 2\} \cap \{2; 3\} = \{2\}$
            \item $A$ et $B$ sont \textit{disjoint} quand 
                $A \cap B = \varnothing$
        \end{itemize}
    \item $A \times B$ et $E^{n}$: produit cartésien. \\
        \Ex \begin{itemize}
            \item Ensembles finis: $$
                \setf{1; 2} \times \setf{3, 4} 
              = \setf{(1; 3), (1; 4), (2; 3), (2; 4)}
            $$
            \item La droite réelle: $\mathbb{R}$,
                le plan réel: $\mathbb{R} \times \mathbb{R} = \mathbb{R}^2$
        \end{itemize}
    
\end{itemize}

\subsection{Quantificateurs.}
\textsc{Attention}: en utilisant ses symboles, nous \textit{posons} une propriété.
Si une démonstration commence par un quantificateur, 
c'est que nous assumons cette propriété vraie, dès le début.

\begin{itemize}
    \item[\Ii] $\forall x$: \enquote{pour tout}, \enquote{quel que soit} $x$. 
        C'est un quantificateur \textit{puissant}:
        pour \textit{tout} ses $x$ possible, les propriétés qui découlent 
        doit être vraies et prouvées. \\
        \Ex 
        \begin{equation}    
            \forall x \in \mathbb{R}, e^x > 0
        \end{equation}
        
    \item[\Ii] $\exists x$: \enquote{il existe} au moins un $x$. \\
        \Ex 
        \begin{equation}
            \forall y \in \mathbb{R}, \quad 
            \exists x \in \mathbb{R}, \quad
            y = x^2
        \end{equation}

    \item $\nexists x$: \enquote{il n'existe aucun} $x$. 
    \item $\exists! x$: \enquote{il existe} un \textit{unique} $x$.
\end{itemize}


\subsection{Logique.}
Bien utiliser les connecteurs et les implications logiques
est crucial mais des contre-sens peuvent être commis.

\paragraph{Symboles.}
\begin{itemize}
    \item[\Ii] $a = b$: $a$ \enquote{est égal} à, est défini par, $b$; 
        défini l'identité entre 2 expressions 
        (scalaires, vecteurs, ensembles, $\dots$).
    \item $a \equiv b \left[m\right]$: 
        $a$ \enquote{est congru} à $b$ modulo $m$ (arithmétique) si, et seulement si,
        $\exists k \in \mathbb{Z} \tq (b - a) = km$ ; 
    \item[\Ii] $\text{A} \Rightarrow  \text{B}$: 
        $\text{A}$ \enquote{donc}, implique, $\text{B}$. Aussi rédigé: \enquote{$
            \underline{Si}    \text{ A } 
            \underline{alors} \text{ B } 
        $}
    \item[\Ii] $\text{A} \Leftrightarrow \text{B}$: 
        $\text{A}$ \enquote{si, et \textit{seulement si}}, réciproquement, 
        $\text{B}$. C'est un relation d'\textit{équivalence} 
        et de \textit{double implication}: $$
        \underline{Si} \text{ A} \Leftrightarrow \text{B } \underline{alors}
        \begin{cases}
            \text{B} \Leftrightarrow \text{A} \\
            \text{A} \Rightarrow  \text{B}    \\
            \text{B} \Rightarrow  \text{A}    \\
        \end{cases} $$
\end{itemize}

\Ex $$
\begin{aligned}
    &x^2 = 25 \\
\Rightarrow \; &x = \sqrt{25} \\
\end{aligned}
$$

Ici, non pas d'équivalence ($\Leftrightarrow$), car cela impliquerait que
$x$ vaille \textit{seulement et uniquement} $\sqrt{25}$,
or $x = \pm\sqrt{25}$

\paragraph{Propriétés.} Posons $\mathcal{P}: $\enquote{
    \underline{si} A \underline{alors} B
}
\begin{itemize}
    \item[\Ii] Si la \textit{réciproque} de $\mathcal{P}$ est vraie, alors:
        \enquote{\underline{$si$} B \underline{$alors$} A}
    \item Si la \textit{contraposée} de $\mathcal{P}$ est vraie:  
        \enquote{\underline{$si$} non B \underline{$alors$} non A }
\end{itemize}

\Ex le théorème de Pythagore: $$
\begin{aligned}
                        &\mathcal{P}: 
        \text{$ABC$ rectangle en $A$ } 
        \underline{donc} \; BC^2 = AB^2 + AC^2 \\
    \text{Réciproque }  &\mathcal{P}: 
        BC^2 = AB^2 + AC^2 \;
        \underline{donc} \text{ $ABC$ rectangle en $A$ } \\
    \text{Contraposée } &\mathcal{P}: 
        BC^2 \neq AB^2 + AC^2 \;
        \underline{donc} \text{ $ABC$ n'est pas rectangle en $A$ } \\
\end{aligned} $$ 

\paragraph{Relations.} Pour une operation $*$ quelconque ($+, -, \times, \dots$)
et des valeurs dans un ensemble $E$.
\begin{itemize}
    \item[\Il] Réflexivité: $x = x$.
    \item[\Il] Symétrie: $x = y \iff y = x$:
    \item[\Il] Transitivité: $
        x = y \text{ et } y = z \;
        \Rightarrow x = z 
    $.
    \item Distributivité avec l'addition: $a * (b + c) = a * b + a * c$
    \item Associativité: $a * (b * c) = (a * b) * c$.
    \item Commutativité: $a * b = b * a$.
    \item Opposé: $\forall x \in E, \exists y \in E
        \tq x * y = 0$
    \item[\Ii] Inverse: $\forall x \in E, \exists y \in E
        \tq x * y = x * x^{-1} = 1 $
    \item[\Ii] Linéarité d'une expression $f$, de 2 constantes $a$, $b$ 
    et d'une inconnue $x$: $f\of{ax + b} = af\of{x} + f\of{x}$
\end{itemize}


\subsection{Divers.}

\paragraph{Nombres.}
\begin{itemize}
    \item $\pi \approx 3.1415926535897932384626433$: Pi, la moitié de la circonférence 
        d'un cercle de rayon $r = 1$.
    \item $e   \approx 2.7182818284590452353602874$: le nombre d'Euler, base de l'exponentielle 
        et du logarithme népérien.
    \item $\floor{x} = \text{E}(x)$: partie entière, \enquote{sol} de $x$. 
        $\floor{x} \leq x < \floor{x} + 1$   \\
        \Ex $$
        \begin{aligned}
            \floor{1.5}         &= 1                \\
            \floor{-4.1}        &= -5               \\
        \end{aligned}
        $$
    \item[\In] $\ceil{x}$: \enquote{plafond} de $x$. $\ceil{x} = \floor{x} + 1$.
    \item $x = \pm y$: \enquote{plus ou moins}, $x \in \setf{+y; -y}$.
    \item[\Ii] $\abs{x}$: valeur absolue de $x$ ou cardinal.
    \begin{itemize}
        \item Pour $x \in \mathbb{R}$, distance à zéro: 
            $\abs{x} = \sqrt{x^2} = \begin{cases}
                \phantom{-}x  &\text{ si } x > 0 \\
                -x  &\text{ sinon }    \\ 
            \end{cases}$ 
        \item Pour $z = a + ib$, distance à zéro: 
            $\abs{z} = \sqrt{a^2 + b^2}$
    \end{itemize}
    \item $\bar{z}$: conjugué de $z \in \mathbb{C}$.
        Pour $z = a + ib,\; \bar{z} = a - ib$
    \item $\arg\of{z}$: argument de $z \in \mathbb{C}$.
        Pour $z = re^{i\theta},\; \arg\of{z} = \theta$
    \item $\ln\of{x}$: logarithme népérien ou naturel de $x \in \mathbb{R}_{+}$.
    \item[\Il] $\log_{b}\of{x}$: logarithme de base $b \in \mathbb{R}_{+}$ 
        de $x \in \mathbb{R}_{+}$ \\
        En particulier $\log_{e}\of{x} = \ln\of{x}$.
\end{itemize}

\paragraph{Fonction et suites.}
Pour un ensemble $E$.
Nous pouvons \textit{étendre} $E$: $$
    (\sim) \quad E_{\infty} = E \cup \setf{+\infty; -\infty}
$$

\begin{itemize}
    \item[\Ii] $a \to b$ pour $a \in E, b \in E_{\infty}$: 
        nous faisons tendre $a$ vers $b$. \\
        \textsc{Attention}: Tout de même, $a \neq b$.
    \item Pour $l \in E_{\infty}$: 
        nous supposons l'existence d'une limite $l$ pour $f(a)$ quand $a$ tends vers $b$: 
        $$
            \lim_{a \to b}{f(a)} = l
        $$

    \item $(u_n)_{n > 0}, \forall n \in \mathbb{N}$: la suite $(u_n)$,
        pour $n > 0$.
    \item[\Ii] Pour des ensembles $A$ et $B$, $f$ est un fonction qui à l'antécédent $x \in A$ 
        applique l'image $f(x) \in B$: $$
        \begin{aligned}
            f:\;&A &\to     \;\; &B \\
                &x &\mapsto \;\; &f(x)
        \end{aligned} $$ 
        \Ex $$
        \begin{aligned}
            f:\;&\mathbb{R}^{\star} &\to \;\; &\mathbb{R}_{+} \\
                &x &\mapsto \;\; &e^{\frac{1}{x}}
        \end{aligned} $$ 
\end{itemize}


\paragraph{Analyse.}
Pour un ensemble $E$.
\begin{itemize}
    \item[\Ii] Somme des valeurs $f(i)$, pour $i$
        allant de $a$ à $b$ \textit{inclus}: $$
            \sum_{i=a}^{b}{f(i)} 
          = f(a) + f(a+1) + \dots + f(b-1) + f(b)
         $$
        \Ex $$
        \begin{aligned}
            &\sum_{i=1}^{3}{2i} 
          = 2\times1 + 2\times2 + 2\times3 \\
          = \; &2\sum_{i=1}^{3}{i}
          = 2 \times (1 + 2 + 3)
          = 12
        \end{aligned}
        $$
    
    \item Produit des valeurs $f(i)$, pour $i$
        allant de $a$ à $b$ \textit{inclus}: $$
            \prod_{i=a}^{b}{f(i)}
          = f(a) \times f(a+1) \times \dots \times f(b-1) \times f(b)
        $$

    \item[\Ii] Dérivée $n$-ième de la fonction $f$ dérivable $n$ fois, pour $d \to 0$, $$
        f^{(n)} \text{ ou } \frac{d^{n}f}{dx^{n}}
    $$
        \Ex $$
            f^\prime = \frac{df}{dx} \quad 
            (f^{\prime})^{\prime} = f^{\prime\prime} = \frac{d^{2}f}{dx^2}
        $$
    \item[\Ii] Pour $a, b \in E$, l'intégrale de $a$ à $b$ de la fonction $f$ 
        $\tq F^{\prime} = f$ pour $d \to 0$,  
        revient à cherche la \textit{primitive} de $f$:
        $$
            \int_{a}^{b}{f\of{x}\,dx} = F\of{x}
        $$
    \item Difference simplifiée: $$
        \brackets{f\of{x}}_{x=a}^{x=b} = f(b) - f(a) 
    $$

        \Ex
        $$
            \int_{0}^{1}{x\,dx}
          = \brackets{\frac{x^2}{2}}_{x=0}^{x=1}
          = \frac{1^2}{2} - \frac{0^2}{2}
          = 0.5
        $$
\end{itemize}

\paragraph{Probabilités.} Pour des événements $A$, $B$ et un variable aléatoire $X$.
\begin{itemize}
    \item $P_{B}\of{A} = \mathbb{P}\brackets{A \vert B}$:
        la probabilité de $A$ \enquote{sachant}, conditionnée par, $B$.
    \item[\Il] $\mathcal{B}\of{n; p}$: loi binomiale 
        de nombre de tirages $n \in \mathbb{N}^{\star}$
        de probabilité $p \in \interval{0}{1}$.
    \item[\Il] $X \sim \mathcal{L}$: $X$ \enquote{suit}, \enquote{s'assimile} à 
        une loi de probabilité $\mathcal{L}$.
    \item $P\of{X = x} = \mathbb{P}\brackets{X = x}$:
        la probabilité que $X$ vaille $x$.
    \item $\mathbb{E}\of{X}$: espérance de $X$.
    \item $\mathbb{V}\of{X}$: variance de $X$.
\end{itemize}


\paragraph{Matrices.}
\begin{itemize}
    \item Définir une matrice $A \in \mathcal{M}_{m, n}(\mathbb{R})$ 
        de taille ($m, n$) par chacun de ses coefficients \\ 
        $\forall i, j \in \mathbb{R}^2, a_{i, j} \in \mathbb{R}$:
        $$
        A = \begin{pmatrix}
            a_{1, 1} & a_{1, 2} & \dots & a_{1, n} \\
            a_{2, 1} & a_{2, 2} & \dots & a_{2, 3} \\
            \dots    & \dots    & \dots & \dots    \\
            a_{m, 1} & a_{m, 2} & \dots & a_{m, n} \\
        \end{pmatrix}
    $$
    \item[\Il] Définir une matrice $A \in \mathcal{M}_{m, n}(\mathbb{R})$ 
        de taille ($m, n$): $$
            A = (a_{i, j}), \; \forall i, j \in \mathbb{R}^2
        $$
    \item[\In] Matrice nulle de taille ($n, n$),
        $\forall i, j \in \iinterval{1}{n}$: $$
    \begin{aligned}
        &\mathbb{O}_{n} = (a_{i, j}) \\
        \text{où } &a_{i, j} = 0 \\
    \end{aligned}
    $$ 

    \item Matrice identité (carrée) de taille $n$, 
        $\forall i, j \in \iinterval{1}{n} $: $$
        I_n = \mathbbm{1}_{n} = \begin{pmatrix}
            1 & 0 & \dots & 0 \\
            0 & 1 & \dots & 0 \\
            \dots & \dots & \dots & \dots \\
            0 & 0 & \dots & 1 \\
        \end{pmatrix}
    $$ 

    \item Déterminant d'un matrice: $$
        det\of{A} = \begin{vmatrix}
            a_{1, 1} & a_{1, 2} & \dots & a_{1, n} \\
            a_{2, 1} & a_{2, 2} & \dots & a_{2, 3} \\
            \dots    & \dots    & \dots & \dots    \\
            a_{m, 1} & a_{m, 2} & \dots & a_{m, n} \\
        \end{vmatrix}
    $$

    
\end{itemize}

\end{document}
